%%%%%%%%%%%%%%%%%%%%%%%%%%%%%%%%%%%%%%%%%%%%%%%%%%%%%%%%%%%%%%%%%%%%%%%%%%%
\chapter{Evaluating Embedding Models: MTEB}
\label{ch:mteb}
%%%%%%%%%%%%%%%%%%%%%%%%%%%%%%%%%%%%%%%%%%%%%%%%%%%%%%%%%%%%%%%%%%%%%%%%%%%

MTEB stands for \textbf{Massive Text Embedding Benchmark}. It is the primary public benchmark
suite for evaluating embedding models across many task families rather than a single retrieval
dataset. That breadth is why it is useful---and also why it is easy to overread.

%%%%%%%%%%%%%%%%%%%%%%%%%%%%%%%%%%%%%%%%%%%%%%%%%%%%%%%%%%%%%%%%%%%%%%%%%%%
\section{What MTEB Measures}
%%%%%%%%%%%%%%%%%%%%%%%%%%%%%%%%%%%%%%%%%%%%%%%%%%%%%%%%%%%%%%%%%%%%%%%%%%%

MTEB evaluates models across a mix of task categories:

\begin{itemize}
    \item Retrieval
    \item Reranking
    \item Classification
    \item Clustering
    \item Semantic Textual Similarity (STS)
    \item Pair classification
    \item Summarization-related matching
    \item Bitext mining
\end{itemize}

The exact benchmark composition depends on which leaderboard slice you are looking at:
English MTEB, multilingual MTEB (MMTEB), or code-focused subsets such as MTEB-Code. A single
leaderboard score is a compressed summary across many task types, domains, and datasets.

%%%%%%%%%%%%%%%%%%%%%%%%%%%%%%%%%%%%%%%%%%%%%%%%%%%%%%%%%%%%%%%%%%%%%%%%%%%
\section{MTEB Task Taxonomy}
%%%%%%%%%%%%%%%%%%%%%%%%%%%%%%%%%%%%%%%%%%%%%%%%%%%%%%%%%%%%%%%%%%%%%%%%%%%

The official MTEB task taxonomy includes these core text task families:

\begin{description}
    \item[\texttt{Retrieval}] Retrieve relevant documents for a query from a corpus.
    \item[\texttt{Reranking}] Reorder a candidate list for better top-of-list relevance.
    \item[\texttt{Classification}] Assign a label to a text.
    \item[\texttt{MultilabelClassification}] Assign multiple labels to a text.
    \item[\texttt{Clustering}] Group texts by semantic similarity without labels.
    \item[\texttt{PairClassification}] Classify the relationship between two texts.
    \item[\texttt{STS}] Semantic textual similarity between two texts.
    \item[\texttt{Summarization}] Evaluate summary-document semantic correspondence.
    \item[\texttt{BitextMining}] Find semantically matching sentence pairs across languages.
    \item[\texttt{InstructionRetrieval}] Retrieval tasks where instruction prompting is part of evaluation.
    \item[\texttt{InstructionReranking}] Reranking tasks where instruction prompting is part of evaluation.
\end{description}

For a local RAG stack, these are the most important by priority:

\begin{enumerate}
    \item \texttt{Retrieval}
    \item \texttt{Reranking}
    \item \texttt{InstructionRetrieval}
    \item \texttt{InstructionReranking}
    \item \texttt{BitextMining} (if multilingual search matters)
\end{enumerate}

Everything else is useful background signal, but it is not the center of gravity for local RAG.

%%%%%%%%%%%%%%%%%%%%%%%%%%%%%%%%%%%%%%%%%%%%%%%%%%%%%%%%%%%%%%%%%%%%%%%%%%%
\section{How to Read an MTEB Score}
%%%%%%%%%%%%%%%%%%%%%%%%%%%%%%%%%%%%%%%%%%%%%%%%%%%%%%%%%%%%%%%%%%%%%%%%%%%

A high MTEB score generally means the model is broadly capable across many embedding tasks.
It does not automatically mean:

\begin{itemize}
    \item It is best for your exact workload
    \item It is fastest
    \item It is the easiest to serve locally
    \item It is best for first-stage retrieval versus reranking
\end{itemize}

For local RAG work, MTEB should be treated as a strong screening benchmark---not the only
decision criterion.

%%%%%%%%%%%%%%%%%%%%%%%%%%%%%%%%%%%%%%%%%%%%%%%%%%%%%%%%%%%%%%%%%%%%%%%%%%%
\section{Common MTEB Metrics}
%%%%%%%%%%%%%%%%%%%%%%%%%%%%%%%%%%%%%%%%%%%%%%%%%%%%%%%%%%%%%%%%%%%%%%%%%%%

MTEB does not use one single metric for every task. It uses task-appropriate metrics, then
aggregates results.

\subsection{What \texttt{@k} Means}

When you see metrics like \texttt{nDCG@10}, \texttt{Recall@20}, or \texttt{Precision@5},
the \texttt{@k} part means: evaluate only the top $k$ ranked items.

\begin{itemize}
    \item \texttt{nDCG@10} --- score the ranking quality of only the top 10 retrieved items
    \item \texttt{Recall@20} --- how many relevant items were captured inside the top 20
    \item \texttt{pass@1} --- did the first attempt succeed?
\end{itemize}

Real systems almost never use the entire ranked list. They typically use top 5, top 10, or
top 20. That is why metrics with \texttt{@k} are often more meaningful than a metric over an
unbounded ranking.

\subsection{nDCG --- Normalized Discounted Cumulative Gain}

nDCG is a ranking metric. Intuitively:

\begin{itemize}
    \item It rewards returning relevant items near the top
    \item It rewards putting the most relevant items earlier in the list
    \item It penalizes good results buried lower in the ranking
\end{itemize}

nDCG is built from two pieces: DCG (Discounted Cumulative Gain) and IDCG (Ideal DCG).
The normalized score is:

\[
\text{nDCG@k} = \frac{\text{DCG@k}}{\text{IDCG@k}}
\]

The DCG formula with graded relevance is:

\[
\text{DCG@k} = \sum_{i=1}^{k} \frac{2^{\text{rel}_i} - 1}{\log_2(i + 1)}
\]

\textbf{Worked example.} Suppose the top-5 retrieved documents have relevance labels
$A = [3, 2, 0, 1, 0]$:

\begin{align*}
\text{DCG}_A@5 &= \frac{7}{1} + \frac{3}{1.585} + \frac{0}{2} + \frac{1}{2.322} + \frac{0}{2.585} \\
               &\approx 7 + 1.893 + 0 + 0.431 + 0 = 9.324
\end{align*}

The ideal ordering $[3, 2, 1, 0, 0]$ gives $\text{IDCG@5} \approx 9.393$.

\[
\text{nDCG}_A@5 = \frac{9.324}{9.393} \approx 0.993 \quad \text{(near-ideal ranking)}
\]

A worse ordering $B = [0, 1, 2, 3, 0]$ of the same documents gives
$\text{nDCG}_B@5 \approx 0.548$---significantly lower, even though both rankings retrieved
the same relevant items. nDCG cares about \textbf{order}, not just membership.

\subsection{MAP --- Mean Average Precision}

MAP measures how well relevant items are retrieved across the ranking. It is sensitive to
whether relevant results appear early, not just whether they appear somewhere. Often used
when there may be multiple relevant documents per query.

\subsection{Recall@k}

Recall@k asks: among all relevant items, how many were captured in the top $k$ results?
This is a coverage metric. It is less concerned with the exact order inside the top $k$ than
nDCG. This matters when your pipeline does first-stage retrieval with a generous top-k
followed by second-stage reranking: high recall at the retrieval stage is often more important
than perfect ranking at that stage.

\subsection{Clustering Metrics}

Clustering tasks often use metrics such as V-measure or NMI (Normalized Mutual Information).
If the embedding space is good, unsupervised clusters should align better with real semantic
categories.

%%%%%%%%%%%%%%%%%%%%%%%%%%%%%%%%%%%%%%%%%%%%%%%%%%%%%%%%%%%%%%%%%%%%%%%%%%%
\section{Which MTEB Metrics Matter Most for RAG}
%%%%%%%%%%%%%%%%%%%%%%%%%%%%%%%%%%%%%%%%%%%%%%%%%%%%%%%%%%%%%%%%%%%%%%%%%%%

If your practical use case is local RAG, the most important categories are:

\begin{itemize}
    \item Retrieval metrics like nDCG and Recall@k
    \item Reranking metrics if you plan to add a reranker
    \item Multilingual retrieval metrics if your corpus or queries span multiple languages
\end{itemize}

For a two-stage retrieval pipeline:

\begin{itemize}
    \item Stage-1 embedding retriever should have strong recall
    \item Stage-2 reranker should improve top-of-list ordering
\end{itemize}

\begin{notebox}[Practical rule for local retrieval]
Use \texttt{Recall@k} to judge whether your first-stage embedding retriever is finding enough
candidates. Use \texttt{nDCG@k} to judge whether the highest-ranked candidates are in the
right order. Use reranking to improve \texttt{nDCG@k} after retrieval if recall is already good.
\end{notebox}

%%%%%%%%%%%%%%%%%%%%%%%%%%%%%%%%%%%%%%%%%%%%%%%%%%%%%%%%%%%%%%%%%%%%%%%%%%%
\section{Qwen3 Series: What It Targets on MTEB}
%%%%%%%%%%%%%%%%%%%%%%%%%%%%%%%%%%%%%%%%%%%%%%%%%%%%%%%%%%%%%%%%%%%%%%%%%%%

The Qwen3 Embedding and Qwen3 Reranker series are designed as a matched retrieval stack
rather than isolated single-purpose models. The official model cards describe the series as
targeting:

\begin{itemize}
    \item Text retrieval
    \item Code retrieval
    \item Text classification
    \item Text clustering
    \item Bitext mining
\end{itemize}

These tasks are related, but not identical. For a local RAG use case, text retrieval and code
retrieval are the most relevant. The series reports strong MTEB scores: Qwen3-Embedding-4B
at 69.45 on multilingual MTEB and 74.60 on MTEB English v2 as of June 2025.

\begin{warningbox}[MTEB Score Interpretation]
High MTEB scores may reflect benchmark proximity rather than out-of-distribution
generalization. Models trained on data that overlaps with MTEB test sets can achieve inflated
public scores. See Chapter~\ref{ch:rteb} for the RTEB benchmark designed to detect this.
\end{warningbox}
